\documentclass{article}
\usepackage{amsmath}
\usepackage{amsfonts}
\usepackage{amssymb}
\usepackage{amsthm}
\usepackage{enumitem}

\newtheorem{assumption}{Assumption}
\newtheorem{definition}{Definition}
\newtheorem{theorem}{Theorem}
\newtheorem{proposition}[theorem]{Proposition}

\begin{document}

\section*{Setup and Notation}

We consider a sample of $n$ i.i.d.\ subjects $(X_i, W_i, Y_i(1), Y_i(0)) \sim \mathbb{P}$, where $W_i \in \{0,1\}$ is a scalar treatment, $X_i \in \mathbb{R}^p$ a covariate vector, and $Y_i(w) \in \mathbb{R}$ potential outcomes.  We observe only $Y_i = Y_i(W_i)$ and write $Z_i = (X_i, W_i, Y_i)$.

We define $e(x) = \mathbb{P}(W_i=1\mid X_i=x)$, $m_w(x) = \mathbb{E}[Y_i\mid W_i=w, X_i=x]$, and $\theta(x) = m_1(x)/m_0(x) - 1$.

The target estimand is
\[
  \theta = \frac{\mu_1}{\mu_0} - 1,
  \qquad
  \mu_w = \mathbb{E}\bigl[\phi(X, P)\,Y(w)\bigr].
\]
The weight $\phi(x, s, t)$ \emph{depends on the unknown distribution $P$} through the scalar $\theta$ and the
conditional effect function $\theta(\cdot)$, where $s$ is a placeholder for the scalar $\theta$ and $t$ is a placeholder for the pointwise value $\theta(x)$; the natural evaluation is $\phi(x,\theta,\theta(x))$.  Partial
derivatives are
\[
  \phi_\theta'(x) := \frac{\partial \phi(x,s,t)}{\partial s}\bigg|_{s=\theta,\,t=\theta(x)}, \qquad
  \phi_{\theta(\cdot)}'(x) := \frac{\partial \phi(x, s, t)}{\partial t}\bigg|_{s=\theta,\,t=\theta(x)}.
\]
That is, $\phi_\theta'(x)$ is the derivative of $\phi$ with respect to the scalar estimand $\theta$, while
$\phi_{\theta(\cdot)}'(x)$ is the derivative with respect to the \emph{pointwise} value $\theta(x)$—treating $\phi$ as a function of the third argument $t$ and evaluating at $t=\theta(x)$.
We treat $e(x)$ and the marginal density $p(x)$ as known, so the only unknown nuisance is the conditional
mean pair $(m_0, m_1)$.


\paragraph{Problem:} We are looking to find a semi-parametric efficiency bound for $\theta$ when the weight $\phi$ is allowed to depend on the unknown $P$ through $\theta$ and $\theta(\cdot)$. We are also looking to find the weights, for which this bound is minimal. Specifically, for any submodel $P_\epsilon$:

\[V(\hat \theta) \geq \left(\frac{d\theta}{d\epsilon}\bigg|_0\right)^2\frac{1}{\mathbb{E}[s(z_i)^2]}.\]

\paragraph{Minimax optimal weights:} We then ask which weight family minimises the worst-case efficiency loss relative to the oracle weight, over $\theta$, over $m_0$, and over local departures from $\theta(\cdot)\equiv\theta$; see Theorem \ref{thm:minimax}.

\begin{assumption}[Known $e(x)$ and $p(x)$]
The propensity score $e(\cdot)$ and the marginal density $p(\cdot)$ are known, so the only unknown nuisance is $(m_0,m_1)$.
\end{assumption}

\begin{assumption}[Unconfoundedness and overlap]
$W \perp\!\!\!\perp (Y(1),Y(0))\mid X$ and $e(x)\in(\eta,1-\eta)$ for some $\eta>0$.
\end{assumption}

Under unconfoundedness we can identify $\theta$ as a functional of the observed data distribution:

\[
  \theta = g(\mu_1,\mu_0) = \frac{\mu_1}{\mu_0} - 1,
  \qquad
  \mu_w = \mathbb{E}\bigl[\phi(X, P)\,Y(w)\bigr].
\]

\begin{assumption}[Homoskedasticity]
\[\mathbb{V}[Y\mid X,W] = \sigma^2\,\mathbb{E}[Y\mid X,W] \]
\end{assumption}

Considering the following parametric submodel:

\[
  \epsilon \mapsto f_{Y\mid X,W}(y \mid x, 1; \epsilon) = f_{Y\mid X,W}(y \mid x, 1)\bigl(1 + \epsilon c_1(y,x)\bigr),
\]
and
\[
  \epsilon \mapsto f_{Y\mid X,W}(y \mid x, 0; \epsilon) = f_{Y\mid X,W}(y \mid x, 0)\bigl(1 + \epsilon c_0(y,x)\bigr).
\]

The score takes the form:
\begin{align}
  s(Z) = W c_1(Y,X) + (1-W) c_0(Y,X)
\end{align}
,with $c_w$ satisfying $\mathbb{E}[c_w\mid X,W=w]=0$ for $w=0,1$.

I do not use homoskedasticity assumption to put a constraint on the tangent space. So, I want the estimate to be efficient in a broader class, which includes any form of variance. However, we chose the weights that give the optimal for a homoskedastic DGP.


\section*{Derivation of the Efficient Influence Function}


Since $\theta = g(\mu_1,\mu_0)$ with $g(a,b)=a/b-1$, the delta method gives


\begin{equation}
  \label{eq:delta}
  \frac{d}{d\epsilon}\theta(P_\epsilon)\bigg|_0 = \frac{1}{\mu_0}\left(\frac{d}{d\epsilon}\mu_1(P_\epsilon)\bigg|_0 - (\theta + 1)\,\frac{d}{d\epsilon}\mu_0(P_\epsilon)\bigg|_0\right).
\end{equation}

Now, we will analyse seperately:
\[
  \mu_0\frac{d}{d\epsilon}\,\theta(P_\epsilon)\bigg|_0
  = \underbrace{\frac{d}{d\epsilon}\,\mathbb{E}_{P_\epsilon}\bigl[\phi(X,P_0)\,\left(\frac{W}{e(X)} - \frac{(\theta+1)(1-W)}{1-e(X)}\right)Y\bigr]\bigg|_0}_{\text{IPW 
  term}} + \]
  \[\underbrace{\mathbb{E}_P\!\left[\frac{d\phi(X,P_\epsilon)}{d\epsilon}\bigg|_0
  \left(\frac{W}{e(X)} - \frac{(\theta+1)(1-W)}{1-e(X)}\right)Y\right]}_{\text{weight adjustment}}
\]

Write $\Psi(Z) := \frac{W}{e(X)} - \frac{(\theta+1)(1-W)}{1-e(X)}$.  In the IPW term the weights do not depend on $\epsilon$, so the derivative acts only on the conditional outcome distribution:

\[\underbrace{\frac{d}{d\epsilon}\,\mathbb{E}_{P_\epsilon}\bigl[\phi(X,P_0)\,\Psi(Z)\,Y\bigr]\bigg|_0}_{\text{IPW term}}
= \mathbb{E}_{P_0}\!\left[\phi(X)\,\Psi(Z)\,Y\cdot s(Z)\right] \]\[
= \mathbb{E}_{P_0}\!\left[\phi(X)\,Y\!\left(\frac{W c_1(Y,X)}{e(X)} - \frac{(\theta+1)(1-W) c_0(Y,X)}{1-e(X)}\right)\right].\]

I now, adjusting the EIF to belong to the tangent space:
\[
  \mathrm{EIF}^{(0)}_\theta(Z)
  := \phi(X)\!\left[\frac{W\bigl(Y-m_1(X)\bigr)}{e(X)}
    - \frac{(\theta+1)(1-W)\bigl(Y-m_0(X)\bigr)}{1-e(X)}\right].
\]


Weight adjustment:
\[
  \underbrace{\mathbb{E}_P\!\left[\frac{d\phi(X,P_\epsilon)}{d\epsilon}\bigg|_0
  \,\Psi(Z)\,Y\right]}_{\text{weight adjustment}}
= \mathbb{E}_{P_0}\!\left[\phi(X)\,\Psi(Z)\,Y\cdot \mathbb{E}_{Z' \sim P_0}\!\left[\frac{d\log \phi(X)}{dP}[Z']\,s(Z')\right]\right].
\]

So, in total:
\[
\mathbb{E}_{P_0}\!\left[\mathrm{EIF}^{(0)}_\theta(Z)\,s(Z)\right]
+ \mathbb{E}_{P_0}\!\left[\phi(X)\,\Psi(Z)\,Y\cdot \mathbb{E}_{Z' \sim P_0}\!\left[\frac{d\log\phi(X)}{dP}[Z']\,s(Z')\right]\right].
\]


\subsection*{$\phi(X,P)=\phi(X,\theta,\theta(\cdot))$}

We now restrict the dependence of the weight on $P$ to the parametrization announced in the setup, $\phi(x) = \phi(x,\theta,\theta(\cdot))$, with derivatives $\phi_\theta'$ and $\phi_{\theta(\cdot)}'$.  Treating $\phi$ as a function of the scalar $\theta$ and the pointwise value $\theta(x)$, the chain rule along the submodel $P_\epsilon$ gives
\begin{equation}
  \label{eq:phi_dot}
  \frac{d}{d\epsilon}\phi(X,P_\epsilon)\bigg|_0
  = \phi_\theta'(X)\,\dot\theta + \phi_{\theta(\cdot)}'(X)\,\dot\theta(X),
\end{equation}
where $\dot\theta := \frac{d\theta(P_\epsilon)}{d\epsilon}|_0$ and $\dot\theta(x) := \frac{d\theta(x;P_\epsilon)}{d\epsilon}|_0$.  Substituting \eqref{eq:phi_dot} into the weight-adjustment expectation splits it into two pieces:
\[
  \text{weight adjustment}
  = \underbrace{\dot\theta\;\mathbb{E}\!\left[\phi_\theta'(X)\,\Psi(Z)\,Y\right]}_{T_\theta}
  + \underbrace{\mathbb{E}\!\left[\phi_{\theta(\cdot)}'(X)\,\dot\theta(X)\,\Psi(Z)\,Y\right]}_{T_{\theta(\cdot)}}.
\]

\paragraph{$T_\theta$}
Conditioning on $X$ and using $\mathbb{E}[WY/e(X)\mid X]=m_1(X)$ and $\mathbb{E}[(1-W)Y/(1-e(X))\mid X]=m_0(X)$,
\[
  \mathbb{E}\!\left[\Psi(Z)\,Y\mid X\right]
  = m_1(X) - (\theta+1)\,m_0(X)
  = m_0(X)\bigl(\theta(X)-\theta\bigr),
\]
so that
\begin{equation}
  \label{eq:T_theta}
  T_\theta = A\,\dot\theta,\qquad
  A := \mathbb{E}\!\left[\phi_\theta'(X)\,m_0(X)\bigl(\theta(X)-\theta\bigr)\right].
\end{equation}

\paragraph{$T_{\theta(\cdot)}$.}
Differentiating $\theta(x) = m_1(x)/m_0(x) - 1$ along the submodel,
\[
  \dot\theta(x) = \frac{\dot m_1(x) - (\theta(x)+1)\,\dot m_0(x)}{m_0(x)}.
\]
The submodel's score $s(Z)=W c_1(Y,X) + (1-W) c_0(Y,X)$ with $\mathbb{E}[c_w\mid X,W=w]=0$ gives
\[
  \dot m_w(x) = \mathbb{E}[(Y-m_w(x))\,c_w(Y,x)\mid X=x,W=w].
\]
We rewrite $\mathbb{E}[h(X)\,\dot m_w(X)]$ in two steps.  \emph{Step 1 (inverse-propensity reweighting).}  For any $g(Y,X)$, $\mathbb{E}[g\mid X,W=1] = \mathbb{E}[W g/e(X)\mid X]$ and analogously for $W=0$, so
\[
  \mathbb{E}[h(X)\,\dot m_1(X)] = \mathbb{E}\!\left[h(X)\,\tfrac{W(Y-m_1(X))}{e(X)}\,c_1(Y,X)\right], \quad
  \mathbb{E}[h(X)\,\dot m_0(X)] = \mathbb{E}\!\left[h(X)\,\tfrac{(1-W)(Y-m_0(X))}{1-e(X)}\,c_0(Y,X)\right].
\]
\emph{Step 2 (score identity).}  Since $W^2=W$ and $W(1-W)=0$, multiplying the score by the indicator collapses it to a single component: $W\,s(Z) = W\,c_1(Y,X)$ and $(1-W)\,s(Z) = (1-W)\,c_0(Y,X)$.  Substituting,
\begin{align*}
  \mathbb{E}[h(X)\,\dot m_1(X)]
  &= \mathbb{E}\!\left[h(X)\,\frac{W\,(Y-m_1(X))}{e(X)}\,s(Z)\right],\\
  \mathbb{E}[h(X)\,\dot m_0(X)]
  &= \mathbb{E}\!\left[h(X)\,\frac{(1-W)(Y-m_0(X))}{1-e(X)}\,s(Z)\right].
\end{align*}
Combining these,
\begin{equation}
  \label{eq:Xi_tilde_identity}
  \mathbb{E}\!\left[h(X)\,m_0(X)\,\dot\theta(X)\right]
  = \mathbb{E}\!\left[h(X)\,\tilde\Xi(Z)\,s(Z)\right],
\end{equation}
where 
\[
  \tilde\Xi(Z)
  := \frac{W\,(Y-m_1(X))}{e(X)} - \frac{(\theta(X)+1)(1-W)(Y-m_0(X))}{1-e(X)}.
\]

Now $T_{\theta(\cdot)} = \mathbb{E}[\phi_{\theta(\cdot)}'(X)\,\dot\theta(X)\,\Psi(Z)\,Y]$, and conditioning on $X$ as in $T_\theta$ replaces $\Psi(Z)Y$ by $m_0(X)(\theta(X)-\theta)$,
\begin{equation}
  \label{eq:T_thetax}
  T_{\theta(\cdot)} = \mathbb{E}\!\left[\phi_{\theta(\cdot)}'(X)\bigl(\theta(X)-\theta\bigr)\,m_0(X)\,\dot\theta(X)\right] = \mathbb{E}\!\left[\phi_{\theta(\cdot)}'(X)\bigl(\theta(X)-\theta\bigr)\,\tilde\Xi(Z)\,s(Z)\right].
\end{equation}

\paragraph{Solving the resulting fixed-point equation.}
Plugging \eqref{eq:T_theta}--\eqref{eq:T_thetax} back,
\[
  \mu_0\,\dot\theta
  = \mathbb{E}\!\left[\mathrm{EIF}^{(0)}_\theta(Z)\,s(Z)\right]
    + A\,\dot\theta
    + \mathbb{E}\!\left[\phi_{\theta(\cdot)}'(X)(\theta(X)-\theta)\,\tilde\Xi(Z)\,s(Z)\right].
\]
The scalar $\dot\theta$ appears on both sides; isolating it,
\[
  (\mu_0 - A)\,\dot\theta
  = \mathbb{E}\!\left[\Bigl(\mathrm{EIF}^{(0)}_\theta(Z) + \phi_{\theta(\cdot)}'(X)(\theta(X)-\theta)\,\tilde\Xi(Z)\Bigr)\,s(Z)\right].
\]
Reading off the influence function gives the closed-form result:
\begin{equation}
  \label{eq:eif_full}
  \boxed{\;
  \mathrm{EIF}_\theta(Z)
  = \frac{1}{C}\Bigl\{\mathrm{EIF}^{(0)}_\theta(Z) + \phi_{\theta(\cdot)}'(X)\bigl(\theta(X)-\theta\bigr)\,\tilde\Xi(Z)\Bigr\},\;}
\end{equation}
with
\begin{equation}
  \label{eq:C}
  C := \mu_0 - A
     = \mu_0 - \mathbb{E}\!\left[\phi_\theta'(X)\,m_0(X)\bigl(\theta(X)-\theta\bigr)\right]
     = \mu_0 - \mathbb{E}\!\left[\phi_\theta'(X)\bigl(m_1(X)-(\theta+1)\,m_0(X)\bigr)\right].
\end{equation}

\begin{assumption}[Regularity]
$m_0(x)>0$ almost surely, $\mu_0\in(0,\infty)$, and $C\neq 0$ in \eqref{eq:C}.
\end{assumption}

\section*{Minimax optimal weights}

Throughout this section $\theta$ denotes the root of
\begin{equation}
  \label{eq:pseudotrue}
  \mathbb{E}\bigl[\phi_s(X)\,m_0(X)\,(\theta(X)-s)\bigr]=0,
  \qquad
  \phi_s(x):=\frac{e(x)(1-e(x))}{1+s\,e(x)},
\end{equation}
which is the defining relation $\theta=\mu_1/\mu_0-1$ evaluated at $\phi=\phi_\theta$.  Set
\begin{equation}
  \label{eq:LambdaQh}
  \Lambda_P(x):=\frac{1+\theta(x)}{e(x)}+\frac{(1+\theta)^2}{1-e(x)},
  \qquad
  \frac{dQ}{dP}:=\frac{\phi_\theta(X)\,m_0(X)}{\mu_0},
  \qquad
  h(x):=\frac{\phi_\theta(x)}{e(x)}=\frac{1-e(x)}{1+\theta\,e(x)},
\end{equation}
and $\tau(x):=\log(1+\theta(x))$, $\tau:=\log(1+\theta)$.

\begin{definition}
$\mathcal{F}$ is the set of measurable $\varphi:\mathbb{R}^p\times(-1,\infty)\to(0,\infty)$ with $\varphi(\cdot,s)\in L^2(m_0\Lambda_P)$.  For $\varphi\in\mathcal{F}$, $\theta_\varphi(P)$ solves $\mathbb{E}[\varphi(X,s)\,m_0(X)(\theta(X)-s)]=0$, and the \emph{implemented weight} is $\vartheta_{\varphi,P}:=\varphi(\cdot,\theta_\varphi(P))$.
\end{definition}

\begin{definition}
For measurable $\vartheta:\mathbb{R}^p\to(0,\infty)$,
\[
  V(\vartheta;P):=\frac{\sigma^2\,\mathbb{E}\bigl[\vartheta^2m_0\Lambda_P\bigr]}{\mathbb{E}[\vartheta m_0]^2},
  \qquad
  R(\varphi;P):=\frac{V(\vartheta_{\varphi,P};P)}{\inf_\vartheta V(\vartheta;P)}.
\]
\end{definition}

\begin{definition}
$\mathcal{N}_\delta:=\bigl\{P:\ \mathbb{E}_Q[(\tau(X)-\tau)^2]\le\delta^2\bigr\}$.
\end{definition}

\begin{theorem}
\label{thm:minimax}
Under Assumptions 1--4:
\begin{enumerate}[label=(\roman*)]
\item $\inf_\vartheta V(\vartheta;P)=\sigma^2/\mathbb{E}[m_0/\Lambda_P]$, attained iff $\vartheta\propto\Lambda_P^{-1}$, and
\[
  R(\varphi;P)
  = \frac{\mathbb{E}\bigl[\vartheta_{\varphi,P}^2m_0\Lambda_P\bigr]\;\mathbb{E}\bigl[m_0/\Lambda_P\bigr]}{\mathbb{E}[\vartheta_{\varphi,P}m_0]^2}
  \;\ge\;1,
\]
with equality iff $\vartheta_{\varphi,P}\propto\Lambda_P^{-1}$.
\item If $\theta(\cdot)\equiv s$ then $\Lambda_P^{-1}=\phi_s/(1+s)$, and
\[
  \inf_{\varphi\in\mathcal{F}}\ \sup_{s\in(-1,\infty)}\ \sup_{m_0}\ R(\varphi;P)=1,
\]
attained iff $\varphi(\cdot,s)\propto\phi_s$ for almost every $s$.
\item $R(\phi_\theta;P)=1+\mathbb{V}_Q\bigl[(\tau(X)-\tau)\,h(X)\bigr]+O(\delta^3)$, and
\begin{equation}
  \label{eq:minimaxvalue}
  \inf_{\varphi\in\mathcal{F}}\ \sup_{s\in(-1,\infty)}\ \sup_{m_0}\ \sup_{P\in\mathcal{N}_\delta}\ R(\varphi;P)
  \;=\;1+\delta^2\Bigl(\operatorname*{ess\,sup}_Q h\Bigr)^2+O(\delta^3),
\end{equation}
attained iff $\varphi(\cdot,s)\propto\phi_s$ for almost every $s$.  Under Assumption 2, $\operatorname*{ess\,sup}_Q h\le(1-\eta)/(1+\theta\eta)\le1$.
\end{enumerate}
\end{theorem}

\begin{proof}
\emph{(i)} For $\vartheta$ a function of $x$ alone, $\phi_\theta'=\phi_{\theta(\cdot)}'=0$, so $A=0$ and $C=\mu_0$ in \eqref{eq:C}, and \eqref{eq:eif_full} reduces to $\mathrm{EIF}_\theta=\mu_0^{-1}\mathrm{EIF}^{(0)}_\theta$.  Squaring and conditioning on $(X,W)$ as in \eqref{eq:Vtheta_generic},
\[
  \mathbb{E}[\mathrm{EIF}_\theta^2]
  = \frac{1}{\mathbb{E}[\vartheta m_0]^2}\,
    \mathbb{E}\!\left[\vartheta^2\!\left(\frac{\mathbb{V}[Y\mid X,1]}{e}+\frac{(1+\theta)^2\,\mathbb{V}[Y\mid X,0]}{1-e}\right)\right],
\]
and Assumption 3 with $m_1(X)=m_0(X)(1+\theta(X))$ makes the bracket $\sigma^2m_0\Lambda_P$, so $\mathbb{E}[\mathrm{EIF}_\theta^2]=V(\vartheta;P)$.  With $g:=m_0\Lambda_P$,
\[
  \mathbb{E}[\vartheta m_0]^2=\mathbb{E}\!\left[\bigl(\vartheta\sqrt g\bigr)\bigl(m_0/\sqrt g\bigr)\right]^2
  \le \mathbb{E}\bigl[\vartheta^2g\bigr]\,\mathbb{E}\bigl[m_0^2/g\bigr],
  \qquad m_0^2/g=m_0/\Lambda_P,
\]
with equality iff $\vartheta\propto m_0/g=\Lambda_P^{-1}$.  Rearranging gives the bound on $V$, and dividing gives the displayed $R$.

\emph{(ii)} If $\theta(\cdot)\equiv s$ then $\theta_\varphi(P)=s$ for every $\varphi$, and
\[
  \Lambda_P=\frac{1+s}{e}+\frac{(1+s)^2}{1-e}
  =\frac{(1+s)\bigl[(1-e)+(1+s)e\bigr]}{e(1-e)}
  =\frac{(1+s)(1+se)}{e(1-e)},
\]
so $\Lambda_P^{-1}=\phi_s/(1+s)$, which does not involve $m_0$.  By (i), $R\ge1$ with equality iff $\varphi(\cdot,s)\propto\phi_s$; since $\varphi(\cdot,s)$ may be chosen separately for each $s$ and the condition is free of $m_0$, the supremum equals $1$ exactly under that condition.

\emph{(iii)} Write $\langle u,v\rangle_P:=\mathbb{E}[uv\,m_0\Lambda_P]$ and $\psi_P:=\Lambda_P^{-1}$, so that $\mathbb{E}[\vartheta^2m_0\Lambda_P]=\|\vartheta\|_P^2$, $\mathbb{E}[m_0/\Lambda_P]=\|\psi_P\|_P^2$ and $\mathbb{E}[\vartheta m_0]=\langle\vartheta,\psi_P\rangle_P$.  Then (i) reads
\begin{equation}
  \label{eq:Rangle}
  R=\frac{\|\vartheta\|_P^2\,\|\psi_P\|_P^2}{\langle\vartheta,\psi_P\rangle_P^2},
\end{equation}
and since $R$ is invariant under $\vartheta\mapsto c\vartheta$ we may write $\vartheta=\psi_P+u$ with $\langle u,\psi_P\rangle_P=0$, giving $\|\vartheta\|_P^2=\|\psi_P\|_P^2+\|u\|_P^2$, $\langle\vartheta,\psi_P\rangle_P=\|\psi_P\|_P^2$ and
\begin{equation}
  \label{eq:Rexact}
  R=1+\frac{\|u\|_P^2}{\|\psi_P\|_P^2}.
\end{equation}

Let $\psi_0:=\Lambda_0^{-1}$ with $\Lambda_0:=\Lambda_P$ at $\theta(\cdot)\equiv\theta$.  Only the first term of $\Lambda_P$ depends on $\theta(\cdot)$, so $\Lambda_P-\Lambda_0=(\theta(X)-\theta)/e$ and
\[
  \psi_0-\psi_P=\frac{\Lambda_P-\Lambda_0}{\Lambda_0\Lambda_P}
  =\frac{\theta(X)-\theta}{e\,\Lambda_0}\cdot\frac{1}{\Lambda_P}.
\]
By \eqref{eq:LambdaQh}, $1/(e\Lambda_0)=h/(1+\theta)$, and $(\theta(x)-\theta)/(1+\theta)=e^{\tau(x)-\tau}-1=(\tau(x)-\tau)+O(\delta^2)$, while $\Lambda_P^{-1}=\psi_0+O(\delta)$.  Hence
\[
  \psi_0-\psi_P=w\,\psi_0+O(\delta^2),\qquad w:=(\tau(X)-\tau)\,h(X).
\]
Since $\psi_0=\phi_\theta/(1+\theta)$, we have $\psi_0^2m_0\Lambda_0=m_0/\Lambda_0=\phi_\theta m_0/(1+\theta)$, so for any $v$,
\[
  \|v\psi_0\|_0^2=\frac{\mu_0}{1+\theta}\,\mathbb{E}_Q[v^2],
  \qquad
  \langle v\psi_0,\psi_0\rangle_0=\frac{\mu_0}{1+\theta}\,\mathbb{E}_Q[v],
  \qquad
  \|\psi_0\|_0^2=\frac{\mu_0}{1+\theta}.
\]
Taking $\vartheta=\psi_0$ in \eqref{eq:Rexact}, $u$ is the component of $\psi_0-\psi_P$ orthogonal to $\psi_P$, and the three displays give
\[
  \frac{\|u\|_P^2}{\|\psi_P\|_P^2}=\mathbb{E}_Q[w^2]-\mathbb{E}_Q[w]^2+O(\delta^3)=\mathbb{V}_Q[w]+O(\delta^3),
\]
which is the first claim of (iii).

For the value, $\mathbb{V}_Q[w]\le\mathbb{E}_Q[(\tau(X)-\tau)^2h^2]\le\delta^2(\operatorname*{ess\,sup}_Qh)^2$.  Taking $\tau(X)-\tau=\delta(\mathbf{1}_A-Q(A))\{Q(A)(1-Q(A))\}^{-1/2}$ with $Q(A)\to0$ and $A$ concentrating where $h$ approaches its essential supremum gives $\mathbb{E}_Q[w]\to0$ and $\mathbb{E}_Q[w^2]\to\delta^2(\operatorname*{ess\,sup}_Qh)^2$, so the bound is approached.

For optimality, by (ii) any $\varphi$ with $\sup R=1+O(\delta^2)$ has $\vartheta_{\varphi,P}=\psi_0+u_0$ with $\langle u_0,\psi_0\rangle_0=0$ and $\|u_0\|_0=O(\delta)$.  Then $\vartheta_{\varphi,P}-\psi_P=u_0+(\psi_0-\psi_P)+O(\delta^2)$, and $\psi_0-\psi_P$ changes sign when $\tau(\cdot)-\tau$ does, while $\mathcal{N}_\delta$ is invariant under that reflection.  Applying $\|a+b\|^2+\|a-b\|^2=2\|a\|^2+2\|b\|^2$ to $a=u_0$, $b=\psi_0-\psi_P$,
\[
  \max_\pm\bigl\|u_0\pm(\psi_0-\psi_P)\bigr\|_0^2\ \ge\ \|u_0\|_0^2+\|\psi_0-\psi_P\|_0^2,
\]
so by \eqref{eq:Rexact} the supremum over $\mathcal{N}_\delta$ is minimised at $u_0=0$, i.e.\ at $\varphi(\cdot,s)\propto\phi_s$.

Finally $\partial_e\{(1-e)/(1+\theta e)\}=-(1+\theta)/(1+\theta e)^2<0$, so $h$ is decreasing in $e$ and Assumption 2 gives $h\le(1-\eta)/(1+\theta\eta)$; and $1+\theta\eta\ge1-\eta$ for $\theta\ge-1$, so this is at most $1$.
\end{proof}

\begin{proposition}
At $\varphi=\phi_\theta$ and $\theta(\cdot)\equiv\theta$,
\[
  V(\phi_\theta;P)=\frac{\sigma^2(1+\theta)}{\mathbb{E}[\phi_\theta(X)m_0(X)]},
  \qquad
  V\bigl(\phi_\theta;P\bigr)\big/(1+\theta)^2=\frac{\sigma^2}{(1+\theta)\,\mathbb{E}[\phi_\theta(X)m_0(X)]}
\]
for $\theta$ and for $\tau=\log(1+\theta)$ respectively.
\end{proposition}

\begin{proof}
By (i) and (ii) of Theorem \ref{thm:minimax}, $V=\sigma^2/\mathbb{E}[m_0/\Lambda_0]$ with $m_0/\Lambda_0=\phi_\theta m_0/(1+\theta)$.  The second expression is the delta-method image of the first under $\theta\mapsto\log(1+\theta)$.
\end{proof}

\section*{Superseded formulation}

\paragraph{The problem.} Define the relative-efficiency loss
\[
  L(\phi;\theta,m_1) \;:=\; \frac{V_\theta(\phi;\theta,m_1)}{V_\theta(\phi^{\text{ref}};\theta,m_1)},
  \qquad \phi^{\text{ref}}(X) := \frac{e(X)(1-e(X))}{1+\theta\,e(X)},
\]
and solve
\begin{equation}
  \label{eq:minimax}
  \inf_{\phi\in\mathcal{F}}\;\sup_{m_1\in\mathcal{T}(\phi)} L(\phi;\theta,m_1).
\end{equation}

\paragraph{The claim.} The saddle point is $(\phi^\star,m_1^\star) = (\phi^{\text{ref}},\,(\theta+1)m_0)$, and the minimax value equals $1$.

\paragraph{Verification.} It suffices to check the two saddle-point inequalities
\begin{equation}
  \label{eq:saddle}
  L(\phi^{\text{ref}};\theta,m_1) \;\le\; L(\phi^{\text{ref}};\theta,m_1^\star) \;\le\; L(\phi;\theta,m_1^\star)
  \qquad\forall\,\phi\in\mathcal{F},\;m_1\in\mathcal{T}(\phi^{\text{ref}}).
\end{equation}

\emph{Left inequality (max in $m_1$):} $L(\phi^{\text{ref}};\theta,m_1)\equiv 1$ by definition, so both sides equal~$1$.

\emph{Right inequality (min in $\phi$):} At $m_1^\star=(\theta+1)m_0$ we have $\theta(X)\equiv\theta$, so the factor $\theta(X)-\theta$ kills both the second term of \eqref{eq:eif_full} and the correction $A$ in \eqref{eq:C}.  Thus $C=\mu_0$ and the influence function collapses to
\[
  \mathrm{EIF}_\theta(Z) \;=\; \mu_0^{-1}\,\mathrm{EIF}^{(0)}_\theta(Z)
  \;=\; \frac{\phi(X)}{\mu_0}\!\left[\frac{W\bigl(Y-m_1(X)\bigr)}{e(X)} \;-\; \frac{(\theta+1)(1-W)\bigl(Y-m_0(X)\bigr)}{1-e(X)}\right].
\]

\emph{Variance computation.} The asymptotic variance is $V_\theta(\phi;\theta,m_1^\star) = \mathbb{E}[\mathrm{EIF}_\theta(Z)^2]$.  Squaring the EIF the cross term vanishes because $W(1-W)=0$, leaving
\[
  \mathrm{EIF}_\theta(Z)^2
  = \frac{\phi(X)^2}{\mu_0^2}\!\left[\frac{W\bigl(Y-m_1(X)\bigr)^2}{e(X)^2}
    + \frac{(\theta+1)^2(1-W)\bigl(Y-m_0(X)\bigr)^2}{(1-e(X))^2}\right].
\]
Conditioning on $(X,W)$ and using $\mathbb{E}[W\mid X]=e$, $\mathbb{E}[1-W\mid X]=1-e$, and $\mathbb{E}[(Y-m_W)^2\mid X,W] = \mathbb{V}[Y\mid X,W]$,
\begin{equation}
  \label{eq:Vtheta_generic}
  V_\theta(\phi;\theta,m_1^\star)
  = \frac{1}{\mu_0^2}\,\mathbb{E}\!\left[\phi(X)^2\!\left(\frac{\mathbb{V}[Y\mid X,1]}{e(X)} + \frac{(\theta+1)^2\,\mathbb{V}[Y\mid X,0]}{1-e(X)}\right)\right].
\end{equation}

\emph{Plugging in homoskedasticity} $\mathbb{V}[Y\mid X,W]=\sigma^2 m_W(X)$, at $m_1^\star=(\theta+1)m_0$ each branch is proportional to $m_0(X)$:
\[
  \frac{\mathbb{V}[Y\mid X,1]}{e} \;=\; \frac{\sigma^2(\theta+1)\,m_0(X)}{e},
  \qquad
  \frac{(\theta+1)^2\,\mathbb{V}[Y\mid X,0]}{1-e} \;=\; \frac{\sigma^2(\theta+1)^2\,m_0(X)}{1-e}.
\]
Factoring $\sigma^2(\theta+1)\,m_0(X)$ out of both summands and putting them over the common denominator $e(1-e)$,
\[
  \frac{\sigma^2(\theta+1)\,m_0(X)}{e} + \frac{\sigma^2(\theta+1)^2\,m_0(X)}{1-e}
  \;=\; \sigma^2(\theta+1)\,m_0(X)\,\frac{(1-e) + (\theta+1)\,e}{e(1-e)}
  \;=\; \sigma^2 m_0(X)\,\frac{(\theta+1)(1+\theta e)}{e(1-e)},
\]
since $(1-e)+(\theta+1)e = 1+\theta e$.  Substituting back into \eqref{eq:Vtheta_generic},
\begin{equation}
  \label{eq:Vtheta_hom}
  V_\theta(\phi;\theta,m_1^\star)
  \;=\; \frac{\sigma^2}{\mu_0(\phi)^2}\,\mathbb{E}\!\left[m_0(X)\,\phi(X)^2\,\Lambda^{(1)}_{\text{hom}}(X)\right],
  \qquad
  \Lambda^{(1)}_{\text{hom}}(X) \;=\; \frac{(\theta+1)(1+\theta e)}{e(1-e)}.
\end{equation}

\emph{Reduction to a Rayleigh quotient.}  Recall $\mu_0(\phi) = \mathbb{E}[\phi(X)\,Y(0)] = \mathbb{E}[\phi(X)\,m_0(X)]$ by iterated expectations, so \eqref{eq:Vtheta_hom} has the form
\[
  V_\theta(\phi;\theta,m_1^\star)
  \;=\; \frac{\sigma^2\,\mathbb{E}\!\left[\phi^2\,g\right]}{\bigl(\mathbb{E}\!\left[\phi\,h\right]\bigr)^2},
  \qquad
  g(X) := m_0(X)\,\Lambda^{(1)}_{\text{hom}}(X),
  \quad
  h(X) := m_0(X).
\]

\emph{Cauchy--Schwarz step.}  Write the inner product
\[
  \mathbb{E}[\phi\,h] \;=\; \mathbb{E}\!\left[\bigl(\phi\sqrt{g}\bigr)\cdot\bigl(h/\sqrt{g}\bigr)\right].
\]
The Cauchy--Schwarz inequality then gives:
\[
  \bigl(\mathbb{E}[\phi\,h]\bigr)^2 \;\le\; \mathbb{E}\!\left[\phi^2\,g\right]\cdot\mathbb{E}\!\left[h^2/g\right],
\],or rearranging:
\[
  \frac{\mathbb{E}[\phi^2 g]}{\bigl(\mathbb{E}[\phi h]\bigr)^2} \;\ge\; \frac{1}{\mathbb{E}[h^2/g]},
\]
with equality iff $\phi\sqrt g \propto h/\sqrt g$:
\[
  \phi(X) \;\propto\; \frac{h(X)}{g(X)} \;=\; \frac{1}{\Lambda^{(1)}_{\text{hom}}(X)} \;=\; \frac{e(1-e)}{(\theta+1)(1+\theta e)} \;\propto\; \phi^{\text{ref}}(X).
\]

Both inequalities in \eqref{eq:saddle} hold, so $(\phi^{\text{ref}},m_1^\star)$ is a saddle point and the minimax value of \eqref{eq:minimax} is~$1$, attained uniquely by $\phi^\star=\phi^{\text{ref}}$.
\end{document}
